\documentclass[10pt]{article}
\usepackage[margin=20mm]{geometry}
\usepackage{amsmath,amssymb,mathtools,fontspec,unicode-math,xurl,hyperref}
\IfFontExistsTF{DejaVu Serif}{\setmainfont{DejaVu Serif}}{\setmainfont{Latin Modern Roman}}
\setmathfont{latinmodern-math.otf}
\setlength{\parindent}{0pt}\setlength{\parskip}{6pt}\sloppy
\title{Technical Verification of
Quasi-invariance of Gaussian measures for the 3d energy critical nonlinear Schrödinger equation}
\author{Initial scientific review: M25-0416}
\date{29 October 2026}
\begin{document}
\maketitle
Abstract: The inspected source is Sun and Tzvetkov’s arXiv:2308.12758v3. The review checks the measure-transport argument, selected normal-form and pairing estimates, and the approximation appendix. One local defect is established: Proposition A.9 does not follow from Definition A.8 and is false with its unrestricted quantifiers. A smooth three-mode initial condition gives an exact counterexample. A uniform bound on the whole orbit restores the refinement assertion; such bounds are available from mass and energy conservation in the subsequent compact-set argument. No failure of the main quasi-invariance theorem is established. The supported assessment is Minor. Independent scientific adjudication remains pending.\par
\section*{Source and scope}
Chenmin Sun and Nikolay Tzvetkov, Quasi-invariance of Gaussian measures for the 3d energy critical nonlinear Schrödinger equation, Communications on Pure and Applied Mathematics 78 (2025), 2305–2353, DOI https://doi.org/10.1002/cpa.70001. The reviewed file is the 49-page author preprint arXiv:2308.12758v3, dated 8 May 2025 [ST]. References to pages below use the PDF page number, which agrees with the printed numbering. Equality of this file with the journal version has not been established.\par
The central statement is equivalence of the Gaussian measure \(\mu_s\) and its image under the defocusing quintic NLS flow on \(\mathbb T^3=\mathbb R^3/(2\pi\mathbb Z)^3\), for every \(s\geq10\) and \(t\in\mathbb R\). The present checks cover the definitions and quantifiers, the finite-dimensional transport and compact-limit argument, selected first- and second-generation estimates, the dangerous cross-generation pairing cancellations, and selected assertions in Appendix A. Text reading and proof verification are distinguished: the review does not independently reconstruct every lattice estimate, every case in Section 8, or the complete external global well-posedness proof. The irrational-torus extension is outside the verified scope.\par
\section*{Definitions relevant to the finding}
Equation reference: eq:nls, eq:energy. Write \(u(t)=\Phi(t)\phi\) for the smooth solution of \[
 i\partial_tu+\Delta u=|u|^4u,\qquad u(0)=\phi.\] Use normalized Haar measure, so that \(\|u\|_{H^1}^2=\sum_{k\in\mathbb Z^3}(1+|k|^2)|\widehat u(k)|^2\). The conserved mass and energy are \[
 M[u]=\|u\|_2^2,\qquad
 H[u]=\tfrac12\|\nabla u\|_2^2+\tfrac16\|u\|_6^6.\] Changing to unnormalized Lebesgue measure multiplies all squared norms and the calculations below by \((2\pi)^3\) and does not change the argument.\par
Equation reference: eq:partition. Definition A.8 on p. 44, paraphrased with its quantifiers retained, calls a finite partition \(-T=\tau_0<\cdots<\tau_m=T\) an \((A,\delta)\)-partition with respect to \(\phi\) if \(A>\delta>0\) and \[
 \|\Phi(\tau_{j-1})\phi\|_{H^1}\leq A,\qquad
 \|\Phi(t)\phi\|_{Z([\tau_{j-1},\tau_j])}<\delta
 \quad(1\leq j\leq m).\] The first bound is imposed at the existing left endpoints, not at every time. The critical \(Z\) norm is defined on p. 38. It decreases upon restriction to a subinterval. On intervals of length less than one it is a sum of two time-\(L^p\) Besov norms, with finite \(p=4.1\) and \(p=100\). Consequently, for a smooth trajectory, \(\|u\|_{Z([-T,T])}\to0\) as \(T\downarrow0\).\par
\section*{A local defect in the refinement assertion}
Problem 1 (An endpoint bound need not survive refinement).  Proposition A.9 (p. 44) asserts that every refinement of an \((A,\delta)\)-partition with respect to \(\phi\) is again an \((A,\delta)\)-partition with the same constants. This statement is false under Definition A.8.\par
\subsection*{Exact counterexample}
Equation reference: eq:phi, eq:coefficients, eq:derivative. Take the smooth initial condition \[
 \phi(x)=1+e^{ix_1}-i e^{2ix_1},\qquad A=\sqrt8,\qquad
 \delta=A/2.\] Its only nonzero Fourier coefficients are \(a_0=1\), \(a_1=1\) and \(a_2=-i\), where the index denotes \((k,0,0)\). Thus \(\|\phi\|_{H^1}^2=1+2+5=8\). Let \(F_k\) be the coefficients of \(|\phi|^4\phi
=\phi^3\overline{\phi}^{\,2}\). Exact multiplication gives \[
 F_0=31+8i,\qquad F_1=31-16i,\qquad F_2=8-31i.\] The remaining coefficients do not enter the derivative at zero, because \(a_k=0\) there. Equation [eq:nls] (\#eq:nls) implies \(a'_k=-i(k^2a_k+F_k)\); hence \[\begin{aligned}
 \left.\frac{d}{dt}\|u(t)\|_{H^1}^2\right|_{t=0}
 &=2\operatorname{Im}\sum_k(1+k^2)\overline{a_k}F_k \notag\\
 &=2\operatorname{Im}\bigl[(31+8i)+2(31-16i)+5i(8-31i)\bigr]
 =32. \end{aligned}\] The calculation has been reproduced by exact Gaussian-integer convolution and a separate SymPy Laurent-polynomial expansion; the retained output gives both methods and every coefficient (Appendix 6 (\#app:evidence)). As consistency checks, \(M[\phi]=3\), \(\|\nabla\phi\|_2^2=5\), \(\|\phi\|_6^6=93\) and \(H[\phi]=18\). At zero the mass derivative is zero, while the kinetic and potential energy derivatives are \(16\) and \(-16\).\par
By smooth local existence, \(f(t)=\|\Phi(t)\phi\|_{H^1}^2\) is continuously differentiable near zero. Equation [eq:derivative] (\#eq:derivative) allows \(T>0\) to be chosen so small that \(f'(t)>0\) throughout \([-T,T]\). Shrink it further so that \(T<1/2\) and \(\|u\|_{Z([-T,T])}<\delta\). Then \[\|u(-T)\|_{H^1}<A=\|u(0)\|_{H^1}
       <\|u(T/2)\|_{H^1}.\] The one-interval partition \(\{[-T,T]\}\) satisfies every condition in [eq:partition] (\#eq:partition). Its refinement at \(T/2\) fails the endpoint condition on the second interval. This is an analytic existence argument for sufficiently small \(T\), not a numerical trajectory or an estimate of a particular positive time.\par
\subsection*{Alternative readings and correction}
The displayed source image confirms that Definition A.8 has no orbit-wide \(H^1\) hypothesis. The small-time construction is admissible: the later restriction \(T\geq1\) in Proposition A.12 is not part of Definition A.8 or Proposition A.9. An implicit conserved \(H^1\) norm cannot resolve the discrepancy; [eq:derivative] (\#eq:derivative) directly excludes it.\par
Equation reference: eq:globalbound. A valid corrected statement is the following: if [eq:partition] (\#eq:partition) holds and, in addition, \[
 \sup_{|t|\leq T}\|\Phi(t)\phi\|_{H^1}\leq A,\] then every refinement is an \((A,\delta)\)-partition. Every new endpoint satisfies [eq:globalbound] (\#eq:globalbound), and restriction monotonicity preserves the strict \(Z\) bound. Equivalently, Definition A.8 can impose the supremum bound on each interval.\par
\subsection*{Impact}
Equation reference: eq:orbit. The common-refinement construction in Proposition A.12, Step 1 (pp. 45–46), cites A.9. This establishes an actual connection to the approximation argument, but does not by itself refute that argument in its application. Conservation and the defocusing sign give \[
 \|\Phi(t)\phi\|_{H^1}^2\leq M[\phi]+2H[\phi].\] For a compact \(K\subset H^1\), the right-hand side has a finite uniform supremum. Choosing the partition constant above its square root gives exactly the additional premise needed for common refinement. On an \(H^1\) ball of radius \(R\), the same bound is a function of \(R\), since \(H^1(\mathbb T^3)\) embeds into \(L^6\). Thus an arbitrary endpoint constant must be replaced by a uniform orbit constant; it need not equal the initial radius.\par
The demonstrated failure concerns the unrestricted constant-preserving assertion A.9. The source’s subsequent proof already invokes energy-based orbit control; the refinement mechanism survives with that premise stated accurately. No failure of the uniform approximation conclusions or Theorem 1.1 has been established. This supports Minor: a local defect without demonstrated central damage. The classification does not follow merely from the appendix location or the brevity of the correction. No author adoption or earlier identification of this particular issue is known from the inspected material.\par
\section*{Checks supporting the main proof structure}
These checks bound the conclusion of the report and record counterevidence against additional objections.\par
\subsection*{Gaussian support and transport}
The series defining \(\mu_s\) has support in \(H^\sigma\) for \(\sigma<s-\tfrac32\), as follows by summing \(\langle k\rangle^{2\sigma}/(1+|k|^{2s})\). The covariance convention using \(\langle k\rangle^{-2s}\) is equivalent: the covariance ratio is \(1+O(|k|^{-2})\), and its squared deviations are summable in three dimensions. For \(F(t)=\rho_{s,N}(\Phi_N(t)A)\), the differential estimate used in Lemma 3.3 (pp. 10–12) has the form \[|F'(t)|\leq C p^\beta F(t)^{1-1/p},\qquad \beta<1.\] It implies \(F(t)^{1/p}\leq F(0)^{1/p}+C|t|p^{\beta-1}\). Sending \(p\) to infinity handles \(F(0)=0\); taking \(p\) comparable to \(\log(1/F(0))\) handles small positive masses and gives \(F(t)\leq C_\varepsilon F(0)^{1-\varepsilon}\). The bounded larger-mass case is absorbed in the constant.\par
Uniform convergence on compact sets, applied also to the inverse flow, gives the inclusion needed on p. 12: \(\Phi(t)K\subset\Phi_N(t)(K+B_\varepsilon)\). Outer regularity passes the estimate to \(K\); inner regularity and a countable union of bounded balls extend it to Borel sets in the Sobolev support. Repeating at \(-t\) gives equivalence of measures. The point-evaluation corollary uses \(\sigma>3/2\) and a nondegenerate Gaussian marginal, and the \(L^1\) isometry uses invertibility and preservation of total variation. These consequences require no extra assertion of a smooth density.\par
\subsection*{Normal forms, pairings and interpolation}
The sign in (2.4) and the correction in (2.8)–(2.10) agree with \(\partial_t e^{-it\Omega}=-i\Omega e^{-it\Omega}\). The resonance cutoff removes the apparent singularity at \(\Omega=0\).\par
For the dangerous opposite-sign cross-generation pair in Section 7, the substitution \[(q_3,q_2,q_5,q_4)=(p_2,p_3,p_4,p_5),\qquad q_1=k_1,\] identifies both constraints and cutoffs. The remaining high-frequency factor is \[\chi_N(k_1)^2|w_{k_2}|^2-
 \chi_N(k_2)^2|w_{k_1}|^2.\] The separated leading term is a real coefficient times a squared modulus, so its imaginary part vanishes. Lemmas 7.4 and 7.5 retain the needed low-frequency gain.\par
An apparent interpolation objection to (7.11) was checked and rejected. Write \(\sigma=s-\tfrac32-\varepsilon\), \(a=1+2\varepsilon\) and \(d=1-2\theta>0\). Even the coarse endpoints \(K^a\) and \(pK^{-d}\) interpolate to \(p^\beta K^{-d/2}\) for \(\beta=1-d/[2(a+d)]<1\), which is summable. Retaining the low-frequency factor \(L_*^2\) gives the stronger endpoint \(pK^{-1}\), because \(L_*^2\prod_j L_j^{3/2-\sigma}\leq1\) when \(\sigma>7/2\). Interpolation with \(\beta=1-2\theta/(a+1)\) then recovers the printed exponent \(-1+2\theta\). This provides counterevidence to treating the compressed interpolation as a defect.\par
Selected Section 8 exponents also admit simultaneous choices. For example, \(\sigma=s-1.51\), \(\delta_0=0.65\), \(\theta=0.32\) and \(\alpha=0.98\) satisfy \(\theta<\delta_0/2<1/3\) and make the checked exponents \(2(s-1-\sigma)-\theta(\sigma-\tfrac72)\) and \(2s+\tfrac32-(2+\alpha)\sigma\) negative for every \(s\geq10\). This verifies these inequalities, not the entire case analysis.\par
\section*{Conclusion}
One substantive local issue is established in the reviewed preprint. The strengthened refinement statement is proved above. The review does not claim that every step of a revised approximation proof has been independently reconstructed, nor that the journal version contains the same issue.\par
Dimension | Recorded state\par
Issue assessment | Minor: unrestricted Proposition A.9 is false.\par
Repair verification | Orbit-wide hypothesis restores refinement exactly; conservation supplies a uniform bound on compact sets.\par
Main conclusion | No refutation of Theorem 1.1 established.\par
Author adoption | Unknown.\par
Scientific adjudication | Pending; no human reviews recorded.\par
Independent acceptance | Not performed in this continuation.\par
\section*{Evidence and limits}
The source PDF has SHA256\par
4503ed3eacaee665bdd30fdab2332842ce
8397218a29d0db28d8958c65b39f96\par
(the line break is typographic only). The source manifest and the complete file bindings are retained under the numeric verification directory. The decisive exact calculation is `checks/approximation/exact\_fourier\_check.json`, with its script, actual execution receipt and standard output. The analytic partition argument and source-page views are in the corresponding `scientific\_check.json`. Primary reading and views are in `checks/root/`; the bounded pairing check is in `checks/cancellation/`. Existing extraction and rendered source bytes were reused, while the recorded scientific checks and indicated image inspections were performed in this continuation.\par
A bounded search of the preserved assigned-paper inventory recovered no substantive predecessor report or mathematical verification. This records absence in the accessible preservation set, not a claim about inaccessible historical chat contents. No completed sibling paper was reviewed. Those roles are not independent acceptance or human review. The controller-supplied product route and runtime self-description are recorded separately in provenance; no backend equivalence is inferred.\par
The complete coverage ledger distinguishes inspected and partially inspected claims, unexamined references, and unavailable journal correspondence. The accompanying build and view records bind the exact current TeX, PDF, retained compiler dependencies and report page images. Their facts are recorded separately from scientific judgment.\par
1 C. Sun and N. Tzvetkov, Quasi-invariance of Gaussian measures for the 3d energy critical nonlinear Schrödinger equation, arXiv:2308.12758v3, 8 May 2025, https://arxiv.org/abs/2308.12758v3. Journal DOI: https://doi.org/10.1002/cpa.70001.\par
\end{document}
